\chapter{Архитектурные развилки и исследовательские продолжения}
\label{ch:tradeoffs}

\section{Зачем обсуждать альтернативы}

Хороший курс не только излагает выбранную архитектуру, но и показывает, от чего она
отказывается. Иначе решения начинают казаться естественными, как будто другой дороги
не было. У архитектуры памяти дорога всегда была.

\section{Граф против плоского векторного хранилища}

\paragraph{Что выигрывает плоское хранилище.}
Дешёвый semantic lookup, простая эксплуатация, понятная масштабируемость.

\paragraph{Что оно проигрывает.}
Явные отношения, адресную ревизию, точное удаление потомков, provenance и
структурный forensic retrieval.

\begin{principle}[Векторный индекс — хороший вход, плохая конституция]
Векторное хранилище полезно как инструмент поиска якорей и похожих эпизодов, но
плохо подходит на роль самой памяти, если память должна поддерживать объяснение,
ретракцию и точную адресность.
\end{principle}

\section{Граф против чисто символической базы}

Символическая система с жёсткими правилами и полноценным TMS превосходит графовую
память в прозрачности и строгости выводов. Но она проигрывает там, где память должна
жить на неполном, шумном и ассоциативном материале.

\paragraph{Выигрыш символики.}
Строгая ретракция, богатая дедукция, точный контроль зависимостей.

\paragraph{Потеря символики.}
Дорогая ручная типизация, слабая терпимость к неполным cue, плохая естественность
ассоциативного вспоминания.

\section{Граф против динамической памяти}

В приложении о динамической памяти уже обсуждаются резервуары, аттракторы и квантовые
ветви. Здесь важно зафиксировать архитектурный смысл сравнения.

\paragraph{Сильная сторона динамики.}
Pattern completion, богатые нелинейные траектории, естественная работа с частичным cue.

\paragraph{Слабая сторона динамики.}
Почти полное разрушение того, ради чего граф строился как юридически и инженерно
адресуемая память: точного forgetting, provenance, явного удаления, объяснимой
ретракции.

\begin{pitfall}[Самый опасный соблазн]
Соблазн состоит в том, чтобы принять красивую динамику за замену всей памяти.
Но значительная часть требований к Рин идёт не из любви к биологическим метафорам,
а из необходимости точной инженерии: кто сказал, когда сказал, что было удалено,
какая гипотеза опровергнута, какие потомки нужно пересчитать.
\end{pitfall}

\section{Новая добавка: правило чистого интерфейса для дальних ветвей}

Это расширение полезно зафиксировать явно.

\begin{principle}[Новая ветвь имеет право жить только за чистым интерфейсом]
Если в систему добавляется резервуарный, аттракторный или квантовый слой, он не
должен получать право произвольно переписывать явный граф. Иначе выигрыши в
ассоциативной динамике будут куплены ценой потери адресности и аудируемости.
\end{principle}

Практически это означает:
\begin{enumerate}
\item динамический слой может предлагать кандидатов;
\item может модулировать retrieval;
\item может служить дополнительным feature-generator;
\item но финальный write-path остаётся у явной памяти.
\end{enumerate}

\section{Почему Stage 1 всё ещё главнее дальних ветвей}

Легко влюбиться в дальнюю ветку раньше, чем построен рабочий позвоночник системы.
Но без Stage~1 невозможно даже понять, что именно новая ветвь улучшает.

\begin{workedbox}[Как выглядит плохой порядок разработки]
Сначала вводится красивая динамическая подсистема. Затем оказывается, что у памяти
до сих пор нет чистого удаления, provenance и нормального Inspector. Тогда любое
улучшение retrieval невозможно честно измерить: неясно, что изменилось из-за новой
динамики, а что ломалось изначально.
\end{workedbox}

\section{Порог входа новой ветви}

Дальняя ветвь заслуживает интеграции только если выполняются три условия.
\begin{enumerate}
\item На базовом стенде сформулирована измеримая слабость текущего графового решения.
\item Новая ветвь атакует именно эту слабость, а не просто выглядит богато.
\item Существует протокол, который сравнивает ветви на одном и том же стенде.
\end{enumerate}

\section{Вывод}

Архитектура памяти не обязана навсегда остаться ровно такой, как в первой рабочей
версии. Но она обязана знать, что в себе является сердцевиной, а что — будущим
расширением. Сердцевина здесь проста: явная адресуемая память, управляемый retrieval,
ревизуемая эпистемика и инженерная проверяемость. Всё остальное должно усиливать
это ядро, а не растворять его.
